\section{反事实比较静态}
\label{sec:counterfactual}

令 $A$ 为已观测的锚环境，$B$ 为目标环境。写为
\begin{equation}
W_J(x,a)=u_J(x,a)+V(\Gamma_J(x,a)),
\qquad J\in\{A,B\},
\label{eq:WJ-decomp-v19}
\end{equation}
其中当前侧差异 $u_B-u_A$ 已知。令
\begin{align}
d_{xa}&:=(u_B-u_A)_{xa},
\label{eq:d-xa-v19}\\
d_{xx}&:=(u_B-u_A)_{xx}.
\label{eq:d-xx-v19}
\end{align}
由于 $H=V\circ\Gamma_A-V\circ\Gamma_B$，第~\ref{sec:shrinking} 节识别的修正项满足
\begin{align}
\boxed{W_{B,xa}=W_{A,xa}+d_{xa}-\eta_{xa},}
\label{eq:transport-xa-v19}\\
\boxed{W_{B,xx}=W_{A,xx}+d_{xx}-\eta_{xx}.}
\label{eq:transport-xx-v19}
\end{align}

令 $C_B(x,a)$ 为目标环境中的货币资源成本。观测到目标基准选择 $x=g_B(a_0)$。假设它是内部严格局部最优点，并且在所考察点上
\begin{equation}
C_{B,xx}-W_{B,xx}>0.
\label{eq:target-regularity-v19}
\end{equation}

\begin{theorem}[反事实比较静态]
\label{thm:counterfactual-v19}
在 $(g_B(a_0),a_0)$ 处，假设环境 $A$ 中的补偿观测识别 $W_{A,xa}$ 与 $W_{A,xx}$，当前侧导数 $d_{xa}$ 与 $d_{xx}$ 已知，并且局部比较设计包含满足式~\eqref{eq:moment-limit}--\eqref{eq:moment-variation}、分别对应 $\partial_{xa}$ 与 $\partial_{xx}$ 的有限差分权重。则 $D_ag_B(a_0)$ 得到点识别，且
\begin{equation}
\boxed{
D_ag_B(a_0)
=
\frac{
W_{A,xa}+d_{xa}-\eta_{xa}-C_{B,xa}
}{
C_{B,xx}-W_{A,xx}-d_{xx}+\eta_{xx}
}.}
\label{eq:main-response-v19}
\end{equation}
\end{theorem}

\begin{proof}
目标环境的一阶条件为
\[
W_{B,x}(g_B(a),a)=C_{B,x}(g_B(a),a).
\]
沿正则局部分支求导，得到
\[
(C_{B,xx}-W_{B,xx})D_ag_B
=
W_{B,xa}-C_{B,xa}.
\]
利用式~\eqref{eq:transport-xa-v19}--\eqref{eq:transport-xx-v19} 以及定理~\ref{thm:primitive-local} 识别 $\eta_{xa}$ 与 $\eta_{xx}$。
\end{proof}

当环境 $A$ 中可以观察邻近补偿比较时，式~\eqref{eq:pi-derivatives} 给出
\begin{equation}
W_{A,xa}=\pi_{A,a},
\qquad
W_{A,xx}=\pi_{A,x},
\label{eq:anchor-pi-derivs-v19}
\end{equation}
于是
\begin{equation}
\boxed{
D_ag_B(a_0)
=
\frac{
\pi_{A,a}+d_{xa}-\eta_{xa}-C_{B,xa}
}{
C_{B,xx}-\pi_{A,x}-d_{xx}+\eta_{xx}
}.}
\label{eq:anchor-to-target-v19}
\end{equation}
当环境 $B$ 中也可以观察邻近补偿比较时，式~\eqref{eq:direct-gprime} 可直接使用。

\subsection{高阶比较静态}
\label{sec:policy-hierarchy}

令
\[
F_B(x,a)=W_{B,x}(x,a)-C_{B,x}(x,a),
\qquad
F_B(g_B(a),a)=0.
\]
重复使用隐式求导是标准做法。一旦 $D_ag_B(a_0),\ldots,D_a^{r-1}g_B(a_0)$ 已知，进入 $D_a^rg_B(a_0)$ 的唯一新的 $(r+1)$ 阶延续项是方向泛函
\begin{equation}
\boxed{
\left(\partial_a+D_ag_B(a_0)\partial_x\right)^r H_x(z_0),
\qquad z_0=(g_B(a_0),a_0),}
\label{eq:response-correction-directional}
\end{equation}
其中微分算子中的系数 $D_ag_B(a_0)$ 保持固定。因此，只要循环矩生成式~\eqref{eq:response-correction-directional}，并结合已知的参考环境导数、当前侧导数和成本导数，定理~\ref{thm:higher-order-differences} 就识别第 $r$ 阶政策导数。无需恢复完整的 $(r+1)$ 阶延续 jet。例如，为识别 $D_a^2g_B(a_0)$ 所需的新三阶延续项为
\[
H_{xaa}+2(D_ag_B)H_{xxa}+(D_ag_B)^2H_{xxx}.
\]
因此，第~\ref{sec:higher-order} 节中的阶次分离直接传递到高阶比较静态。
